\chapter{Overview}
\label{chap:overview}

This chapter is a map, not a summary.  Its purpose is to let a reader who
opens the blueprint for the first time see, in a few minutes, what the VR
Cycle is, what its modules are for, how they depend on one another,
and which chapter to read next.  It states no theorems and cites no Lean
objects; each later chapter does that for its own material.

%──────────────────────────────────────────────────────────────────────
\section{What VR is}
\label{sec:overview_what}

VR is an operational reconstruction of mathematics from a single
base operation --- the empty object~$\emptyset$ --- through a small
fixed repertoire of operations: succession~($t$) in the formal system, and
reference~($\in$) once sets enter.  (Up to v1.0.1 the formal system also
listed implication~($\to$) on a two-valued type as a primitive; it was
removed in v1.1.1 --- it generated only a finite truth-function algebra used
by no VR theorem, and the logic VR reasons with is metatheoretic, not
generated by VR.)  Every
mathematical object in the cycle is built as an operation on~$\emptyset$;
no further primitives are admitted.  The discipline is uniform across all
the works: an object is introduced together with an explicit account of
how it is operationally constructed and what logical axioms its
construction depends upon.

From this root the cycle climbs through the natural numbers, the integers,
rationals, reals and complex numbers, the ZF and ZFA set universes, a
two-register apparatus separating operational from formal reasoning,
operational algebra over the standard structures, an operational
Hahn--Banach theorem via Riesz representation, and a point-free binary
Tychonoff theorem.  Each step is formalised in Lean~4 and carries a
documented axiom profile, so that the logical cost of every object is
visible rather than implicit.

Since 2026-09-12 the repository builds two Lake libraries, and the
relation to mathlib differs between them.  The library \texttt{VR} --- VR
proper: the formal system, the witnessed numbers
$\mathbb{Z}\to\mathbb{Q}\to\mathbb{R}$ on their own pairs, the
operational set universe \texttt{OpSet} (Chapter~\ref{chap:setsop}), the
ZTL layer, the forms, the apparatus, the topology tower up to Tychonoff and
the Brouwer continuum --- imports \emph{no} mathlib at all (the few
borrowed pieces are supplied by its own \texttt{VR/Prelude/}), and every one
of its declarations stands on the empty axiom list, enforced at build time
by \texttt{VR/Guard.lean}.  The library \texttt{VRClassical} --- the
classical register and the bridges --- is built \emph{on} mathlib, not
parallel to it: there the ZFC-reading set type \texttt{OSet} is mathlib's
well-founded \texttt{ZFSet} used directly (\texttt{abbrev OSet := ZFSet}
in \texttt{VRClassical/Sets/Foundation.lean}); the anti-foundational
universe \texttt{OSetZFA} is constructed from mathlib's \texttt{PFunctor.M}
coinductive machinery; the operational Hilbert space wraps mathlib's
classical \texttt{InnerProductSpace}; and the topology bridge instantiates
mathlib's \texttt{Order.Frame}.  \texttt{VRClassical} depends on
\texttt{VR}; nothing in \texttt{VR} depends on it.  The
contribution of the cycle is the coherent synthesis erected above this
foundation, from one root through arithmetic, sets, apparatus and
applications --- not a replacement of the foundation, and not any single
hidden capability of one module.

%──────────────────────────────────────────────────────────────────────
\section{The three architectural layers}
\label{sec:overview_layers}

The nine original modules organise into three layers; the later
chapters are placed after them.

\textbf{Layer~I --- Foundation} (Chapters~\ref{chap:vr}, \ref{chap:numbers},
\ref{chap:sets} and~\ref{chap:setszfa}).  The vertical from~$\emptyset$
upward: the VR formal system fixes the operational base and recovers Peano
arithmetic; VR-Numbers builds the number systems on it; VR-Sets realises
the ZFC universe through \texttt{OSet} (in \texttt{VRClassical}; the
choice-free operational universe of the core is \texttt{OpSet},
Chapter~\ref{chap:setsop}); and VR-Sets-ZFA realises the
anti-foundational universe through \texttt{OSetZFA}.  The two set-theoretic
universes are connected by the \texttt{embedOSet} bridge, so that the
well-founded and non-well-founded worlds sit in one development rather than
two.

\textbf{Layer~II --- Methodology} (Chapters~\ref{chap:forms}
and~\ref{chap:apparatus}).  VR-Forms introduces the two-register apparatus
that distinguishes what is operationally constructed from what is formally
asserted.  VR-Apparatus generalises that apparatus into a reusable
framework: the two operationality modes (\texttt{PredicateOperationality},
\texttt{ReferenceOperationality}), the two transit modes (Mode~A and
Mode~B), and \texttt{InterApparatusMorphism} connecting them.  This layer
is where the discipline used implicitly across the foundation is made
explicit and portable.

\textbf{Layer~III --- Applications} (Chapters~\ref{chap:audit},
\ref{chap:algebra} and~\ref{chap:topology}).  VR-Audit applies the
apparatus to functional analysis, deriving an operational Hahn--Banach
theorem from Riesz representation and exhibiting Mode~B on a classical
analytic result.  VR-Algebra reuses the apparatus across additive groups,
rings, fields and modules, exercising both Mode~A and Mode~B in algebra.
VR-Topology transfers the apparatus to an unfamiliar domain, proving a
choice-free point-free binary Tychonoff.

The dependency direction is strict: each Layer~III application uses
methodology from Layer~II, which in turn rests on the foundation in
Layer~I.  No arrow runs the other way.  This is the \texttt{\textbackslash
uses} chain made visible in the dependency graph.

\textbf{Later chapters.}  VR-Transit (Chapter~\ref{chap:transit}) extends
the apparatus with transit conservativity and a bounded witness library.
Brouwer's fixed-point theorem via Sperner (Chapter~\ref{chap:brouwer}) is
mathlib-bound and not part of the cycle.  The operational continuum and
number spectrum (Chapter~\ref{chap:continuum}), the operational set
universe under Brouwer (Chapter~\ref{chap:setsop}) and the operational sets
under ZTL verdicts of VR Part~II (Chapter~\ref{chap:setsztl}) live in the
core library \texttt{VR}, on the empty axiom list.

%──────────────────────────────────────────────────────────────────────
\section{Dependency structure}
\label{sec:overview_deps}

The dependency graph is methodological, not chronological.  VR-Audit was
completed and published before VR-Apparatus, yet it depends on the
apparatus as the canonical Mode~B instance; the blueprint records the
methodological dependency because that is what a reader needs in order to
navigate the work, not the order in which the modules happened to be
written.  Reading the graph as a timeline would therefore mislead; it is a
map of logical support.

The full dependency graph is rendered separately as the
\href{dep_graph_document.html}{dependency graph} accompanying this
blueprint, where each node links to the corresponding chapter and the
chapter-to-chapter edges are shown visually.

%──────────────────────────────────────────────────────────────────────
\section{Axiom profile across the cycle}
\label{sec:overview_axioms}

Every public object in the cycle has been audited with \texttt{\#print
axioms}, and the results fall into four tiers.  Since 2026-09-12 the tiers
above the first occur only in \texttt{VRClassical}: the whole core library
\texttt{VR} sits in the first.

\begin{itemize}
\item \texttt{[]} --- axiom-free.  The entire core library \texttt{VR}
sits here --- the formal system (Chapter~\ref{chap:vr}), the witnessed
numbers, \texttt{OpSet}, the forms, the apparatus core, the topology tower
including \texttt{tychonoff\_binary}, the continuum and the ZTL layer ---
enforced at build time: \texttt{VR/Guard.lean} runs
\texttt{\#assert\_axiom\_free\_library} over every constant of the
library, and \texttt{lake build VR} fails on any axiom (the only exemption
is the tool code of the tactic and elaborator metaprograms, named in the
guard's output).  Nothing beyond Lean's definitional framework is used.

\item \texttt{[Quot.sound]} --- quotient soundness only.  A sub-ceiling
first isolated in VR-Apparatus (Finding~S2-B), below the usual
constructive ceiling.

\item \texttt{[propext, Quot.sound]} --- propositional extensionality with
quotient soundness.  This was the cycle's standard constructive ceiling
until the 2026-09-12 integrity pass, which brought the core down to
\texttt{[]}; it remains in \texttt{VRClassical} (the quotient numbers, the
Frame bridge).  VR-Topology's \texttt{tychonoff\_binary} was released at
this tier (v1.0.0) and now stands on \texttt{[]}, even though the classical
Tychonoff theorem is equivalent to the Axiom of Choice; the point-free
formulation is what makes the choice-free profile attainable.

\item \texttt{[propext, Classical.choice, Quot.sound]} --- the full mathlib
ceiling.  It enters only through black-box dependence on mathlib's
classical theorems: Riesz representation in VR-Audit, and the ordinal,
foundation and choice mechanisms behind \texttt{ZFSet} in VR-Sets.
\end{itemize}

The methodological point is not that the cycle avoids
\texttt{Classical.choice} as an end in itself.  It is that the cycle
documents precisely where choice enters and structures the apparatus so
that classical machinery lives in the formal register without contaminating
the operational register.  This is the two-register discipline of VR-Forms
applied to the cycle's own logical accounting.

%──────────────────────────────────────────────────────────────────────
\section{What to read first}
\label{sec:overview_reading}

Different readers will want different entry points.

A reader interested in foundations and axiomatic minimalism should start
with the VR formal system (Chapter~\ref{chap:vr}), which is self-contained,
axiom-free, and lays out the cycle's operational position from~$\emptyset$.
VR-Sets (Chapter~\ref{chap:sets}) then shows how that position extends to set
theory through \texttt{OSet}, and Chapter~\ref{chap:setsop} rebuilds it
choice-free as \texttt{OpSet}.

A reader interested in constructive analysis and proof mining should
compare VR-Audit (Chapter~\ref{chap:audit}) with VR-Topology
(Chapter~\ref{chap:topology}): two contrasting analytic settings, one reaching
its result through Mode~B over a black-box Riesz representation, the other
through a point-free formulation that avoids choice entirely.

A reader interested in formalisation methodology should read VR-Apparatus
(Chapter~\ref{chap:apparatus}), the meta-work that formalises the framework
used implicitly throughout the cycle, and then VR-Algebra
(Chapter~\ref{chap:algebra}), which confirms the framework's transferability
across algebraic structures.

A reader interested in non-well-founded set theory should turn to
VR-Sets-ZFA (Chapter~\ref{chap:setszfa}), which formalises the Anti-Foundation
Axiom as a theorem, obtained from the final-coalgebra property of
\texttt{CoPSet} via mathlib's \texttt{PFunctor.M}; the core's
\texttt{OpSet.afa} (Chapter~\ref{chap:setsop}) obtains AFA on the empty
axiom list without that machinery.

%──────────────────────────────────────────────────────────────────────
\section{Position relative to existing programmes}
\label{sec:overview_position}

VR is a synthesis, not a discovery of new mathematics.  Each component
piece has an identifiable predecessor, and the blueprint acknowledges them
plainly.

\begin{itemize}
\item \textbf{Hilbert.}  The formal--ideal distinction is reframed by the
two-register apparatus (Chapter~\ref{chap:forms}), with the conservativity
relation between the registers made a machine-verified claim rather than an
informal one.

\item \textbf{Bishop.}  The cycle shares Bishop's constructive discipline
but differs in strategy: where Bishop rejects classical reasoning, VR
accommodates it inside the formal register and keeps it out of the
operational one.

\item \textbf{Aczel.}  The Anti-Foundation Axiom is Aczel's; VR-Sets-ZFA
(Chapter~\ref{chap:setszfa}) contributes its formalisation as a theorem in
Lean using coinductive types, not the axiom itself.

\item \textbf{Coquand, Vickers and Sambin.}  VR-Topology
(Chapter~\ref{chap:topology}) uses their formal-topology framework directly;
the chapter is explicit that the contribution is transferability of the
apparatus, not a result within the point-free programme.

\item \textbf{Kohlenbach.}  Mode~B's extraction of operational content from
a classical theorem is structurally close in spirit to proof mining, though
the apparatus is distinct from the Dialectica interpretation.
\end{itemize}

What is new is none of these pieces individually.  It is the coherent,
machine-verified synthesis that runs from a single root through arithmetic,
two set universes, the apparatus, and three applications --- one organism
rather than a set of separate notes.  Making that connectivity visible, before
the reader looks at any single chapter, is the work of this overview.

%──────────────────────────────────────────────────────────────────────
\section{Provenance and reproducibility}
\label{sec:overview_provenance}

The Lean~4 formalisation lives at
\href{https://github.com/inventor1975/VRCycle}{github.com/inventor1975/VRCycle}.
The toolchain is pinned in the repository's \texttt{lean-toolchain} file;
\texttt{lake build} compiles both libraries (with linter warnings, no
errors), \texttt{lake build VR} runs the axiom guard over the core, and
the axiom profiles quoted throughout the blueprint are produced by
\texttt{\#print axioms} run on every public object.

Archival records are deposited on Zenodo: ten preprints (VR-Forms
through v1.0.3, VR Part~II and \emph{Choice as an Act} included) and seven
Lean software releases; the repository README lists them with their DOIs.
Some modules are cited differently by deliberate decision --- VR-Algebra by
GitHub tag without a separate Lean DOI, a workflow simplification, and
VR-Topology and VR-Transit by git tag, with no Zenodo record.

The original cycle was developed between 15 and 28~May~2026 using AI assistance in
the Variant~A pattern: a human curator directing an architect model
(Claude Opus~4.7) that produced specifications, and an executor model
(Claude Sonnet~4.6) that implemented them in Lean.  Later passes
(June--October~2026) were made with later Claude models (Opus~4.8,
Fable~5, Fable~5.1, Opus~5.5), recorded in the commits' co-author trailers.  The repository
README carries the full account of authoring; it is recorded here as a
plain fact about how the cycle was produced, not as a framing of the
mathematics, which stands or falls on the machine-checked proofs.
